% XVI-P p18 note: the D11 matching premise M2, D11, and the full hh slot
\section{P.5 calculation, step 4 (p18): M2 discharged, D11 evaluated, the full hh slot}
\scopebox{Scope.}{p18 keeps the p14r contract and the p15, p16 and p17 calculation layers byte-identical
(\texttt{release\_p18.py} recomputes their sha256). It discharges the D11 matching premise M2 on the 4D components,
evaluates D11 under the declared prescription, and rechecks the ch-projected hh slot with D11 included.
quantum\_inputs\_evaluated is earned (every graph of the inventory is evaluated in one prescription).
projected\_identity\_verified, projected\_restoration\_verified and bulk\_symmetry\_restored are NOT earned.
O.1 stays NOT\_CLOSED; the insertion stays frozen; XV-B and XIV-J are unopened; Route B is not discharged; no
$c\,\delta$ is added to $u^R_\mu$.}

\subsection{M2 (certificate \texttt{P18\_D11\_M2.json})}
M2 states that on $U=\mathbb R^4\setminus\{0\}$ the inverse Fourier transform of $P^{DR}_{AB}$ equals
$R_AR_B\,\Pi^R_{\rm tr}$, with $\Pi^R_{\rm tr}=(i\kappa^2/4)\,K u^R_\mu$, $K=\alpha D_2+\beta D_0$,
$\alpha=A^2/60$, $\beta=A^2/18$ and $A=1/4\pi^2$.
\begin{itemize}
\item \emph{Route B (exact; primary).} Write $P^{DR}=\kappa^2N_\varepsilon\hat G$ with
  $\hat G=\sum_sA_s(d)S_s\,{\rm MB}$ and ${\rm MB}=1/\varepsilon+2-L$. At $d\to4$ its nonlocal part is
  $-(i\kappa^2/16\pi^2)L\sum_sA_s(4)S_s(p)$; everything else is polynomial. On the other side,
  $\mathcal F[u^R]=c_{\rm FT}\log(-p^2-i0)+{}$const, and $K\to\hat K(p)$ with
  $\hat P_{mn}=p_mp_n-\eta_{mn}p^2$. This operator is fourth order, so the sign convention of the transform drops out.
  M2 then holds iff, for each of the 100 component pairs,
  $X_{AB}=\sum_sA_s(4)S_s|_{AB}=-4\pi^2c_{\rm FT}A^2\,Y_{AB}=-(i/4)\,Y_{AB}$, where
  $Y=R_AR_B(\hat D_2/60+\hat D_0/18)$. This holds exactly on all 100 pairs. The difference then has a polynomial
  Fourier transform, so it is supported at $\{0\}$, and $\mu$ enters only through polynomials.
\item \emph{$c_{\rm FT}$, calibrated twice.} (a) From the library DR bubble,
  $\mathcal F[D_F^2]=-N_\varepsilon{\rm MB}$ and $D_F^2=A^2u$. (b) As $-r$, with the XVI-O residue $r=-i\pi^2$. Both
  give $c_{\rm FT}=i\pi^2$.
\item \emph{Route A (structural corroboration).} M1 is re-run: the P integrand equals the Fourier convolution of the
  XVI-O product $T=VV(CC+CC)$ (100/100). The frozen XVI-O code, imported read-only, re-derives $\alpha,\beta$ and
  verifies $T=Ku$ on $U$ for all 256 components. A regularised convolution agrees with the pointwise product on $U$
  (cited: XVI-O lem:BV\_rules\_on\_U, Constructed).
\item \emph{Controls and mutations.} D11 is transverse exactly in $d$; the FP and PV loop routes agree. Six mutations
  are each rejected: $\alpha\leftrightarrow\beta$, off-diagonal weight 1, flipped FT sign, doubled symmetry factor,
  D01 in place of D11, and $\beta=0$.
\item \emph{Scope.} The flag is qualified as D11\_matching\_verified\_scope = ``4D components''. M2 is shown on the
  4D components only. The evanescent components are the 8 $\tilde g$-carrying structures, i.e.\ components with an index
  pair in the $\varepsilon$-scalar directions. Their status is not\_applicable\_not\_verified, because XVI-O is
  strictly 4D and has no counterpart. The evanescent structures of D11 have nonzero
  $L$-coefficients (8, recorded); they are $\varepsilon$-scalar components outside the 4D statement.
\item \emph{Consequence (M3/M4).} M3 is now unconditional on the 4D components: $L_{\rm match}$ is supported at 0. It is transverse and homogeneous of degree four,
  so it lies in the 2-dimensional transverse h$\cdot$h subspace. Its coordinates need the XVI-O finite constant of
  $u^R$ and are not computed.
\end{itemize}

\subsection{D11 (certificate \texttt{P18\_D11.json}) and the full slot (\texttt{P18\_hh\_slot.json})}
\begin{itemize}
\item \emph{Prescription.} Operational DRED R1--R6 with opaque $d$, MS-bar, Feynman $+i0$, N4 contact exclusions,
  M3 bounds 4/8/4, exactly-degree-four normalisation. Both points are checked per structure: every structure has
  momentum degree 4, so the order-0/2 contacts excluded by N4 do not arise.
\item \emph{Inventory.} quantum\_inputs\_evaluated is earned. Every graph of the three surviving ch rows is evaluated,
  and none is Missing: hh D01, D02, D03, D11, D12; K$\cdot$c D13; one-point D14a--c.
\item \emph{Evaluation.} 14 structures, 6 of them 4D. The covariant form is exact and unique; the FP and PV routes
  agree exactly in $d$; the result is exchange symmetric. The nonlocal coefficient equals minus the pole. Premature
  $d=4$ is refused and would shift the local parts. Each of the four mutations is rejected.
\item \emph{4D pole.} $z_{D11}=i(-1/160,-1/480,1/80,1/240,-1/120)=i(\tfrac1{240}k_1-\tfrac1{120}k_2)$, which is
  transverse.
\item \emph{Full row.} $p_m G_{D11}^{(mn)(rs)}=0$ holds identically in $d$, for the 4D and the evanescent parts.
  D11 therefore decouples from the slot by itself; D01 and D03 decouple only as a sum.
  \begin{itemize}
  \item $X_{D01+D03+D11}=0$ for all seven external assignments. This includes $\hat c,\tilde h$ (H2), where
    $X_{D01}=-X_{D03}=p^4\,\tilde a\!\cdot\!\tilde b\;\hat c\!\cdot\!p/16$.
  \item No evanescent structure survives, so D11 does not reopen the evanescent question and no counterterm
    revision is needed.
  \item The summed 4D pole is $i(\tfrac{41}{80}k_1-\tfrac{43}{120}k_2)$, which is transverse.
  \item Caveat: D11 phase and symmetry-factor mutations are invisible to the slot. They are tested by M2 instead.
  \item The summed evanescent poles of the hh Green function (8 structures) are nonzero but transverse.
    Renormalising that Green function itself is outside the projection and not claimed.
  \end{itemize}
\end{itemize}
% generated by code/p5_calc/release_p18.py -- do not edit
\noindent M2: \texttt{DISCHARGED (4D components)}; route B: 100 component pairs, 0 mismatches, constant \texttt{-I/4}; c\_FT = \texttt{I*pi\^{}2} (DR bubble) = \texttt{I*pi\^{}2} (-r).
\begin{center}\scriptsize\begin{tabular}{p{3.0cm}p{0.6cm}p{1.5cm}p{1.5cm}p{3.6cm}p{2.4cm}}\hline
D11 structure & evan. & UV pole & coeff. of L & local (MS-bar) & shift (d=4 first) \\\hline
\texttt{g4\_mn g4\_rs (p.p)\^{}2} & no & \texttt{-I/80} & \texttt{I/80} & \texttt{-I/75} & \texttt{7*I/600} \\
\texttt{g4\_mn gt\_rs (p.p)\^{}2} & yes & \texttt{-I/120} & \texttt{I/120} & \texttt{-77*I/3600} & \texttt{-17*I/3600} \\
\texttt{g4\_mn p\_r p\_s (p.p)\^{}1} & no & \texttt{I/80} & \texttt{-I/80} & \texttt{I/75} & \texttt{-7*I/600} \\
\texttt{g4\_mr g4\_ns (p.p)\^{}2} & no & \texttt{-I/480} & \texttt{I/480} & \texttt{-23*I/3600} & \texttt{-I/450} \\
\texttt{g4\_mr gt\_ns (p.p)\^{}2} & yes & \texttt{I/480} & \texttt{-I/480} & \texttt{23*I/3600} & \texttt{I/450} \\
\texttt{g4\_mr p\_n p\_s (p.p)\^{}1} & no & \texttt{I/480} & \texttt{-I/480} & \texttt{23*I/3600} & \texttt{I/450} \\
\texttt{g4\_rs gt\_mn (p.p)\^{}2} & yes & \texttt{-I/120} & \texttt{I/120} & \texttt{-77*I/3600} & \texttt{-17*I/3600} \\
\texttt{g4\_rs p\_m p\_n (p.p)\^{}1} & no & \texttt{I/80} & \texttt{-I/80} & \texttt{I/75} & \texttt{-7*I/600} \\
\texttt{gt\_mn gt\_rs (p.p)\^{}2} & yes & \texttt{-I/480} & \texttt{I/480} & \texttt{-23*I/3600} & \texttt{-I/450} \\
\texttt{gt\_mn p\_r p\_s (p.p)\^{}1} & yes & \texttt{I/120} & \texttt{-I/120} & \texttt{77*I/3600} & \texttt{17*I/3600} \\
\texttt{gt\_mr gt\_ns (p.p)\^{}2} & yes & \texttt{-I/480} & \texttt{I/480} & \texttt{-23*I/3600} & \texttt{-I/450} \\
\texttt{gt\_mr p\_n p\_s (p.p)\^{}1} & yes & \texttt{-I/480} & \texttt{I/480} & \texttt{-23*I/3600} & \texttt{-I/450} \\
\texttt{gt\_rs p\_m p\_n (p.p)\^{}1} & yes & \texttt{I/120} & \texttt{-I/120} & \texttt{77*I/3600} & \texttt{17*I/3600} \\
\texttt{p\_m p\_n p\_r p\_s} & no & \texttt{-I/60} & \texttt{I/60} & \texttt{-47*I/1800} & \texttt{13*I/1800} \\
\hline\end{tabular}\end{center}
\begin{center}\scriptsize\begin{tabular}{p{2.0cm}p{1.3cm}p{1.3cm}p{1.3cm}p{1.3cm}p{1.3cm}p{1.2cm}}\hline
4D pole & $z_1$ & $z_2$ & $z_3$ & $z_4$ & $z_5$ & transverse \\\hline
D11 & \texttt{-I/160} & \texttt{-I/480} & \texttt{I/80} & \texttt{I/240} & \texttt{-I/120} & yes \\
D01+D03 & \texttt{-23*I/240} & \texttt{-61*I/240} & \texttt{23*I/120} & \texttt{61*I/120} & \texttt{-7*I/20} & yes \\
D01+D03+D11 & \texttt{-49*I/480} & \texttt{-41*I/160} & \texttt{49*I/240} & \texttt{41*I/80} & \texttt{-43*I/120} & yes \\
\hline\end{tabular}\end{center}
\begin{center}\scriptsize\begin{tabular}{p{2.6cm}p{2.6cm}p{2.6cm}p{1.2cm}p{1.6cm}}\hline
assignment $(c,h)$ & $X_{D01}$ & $X_{D03}$ & $X_{D11}$ & sum \\\hline
\texttt{Q4S (declared, R2)} & \texttt{p.p\^{}2*(d*a.b*c\^{}.p - 3*a.b*c\^{}.p + a.c*b\^{}.p + a\^{}.p*b.c - a\^{}.p*b~.c~ + 2*a~.b~*c\^{}.p - a~.c~*b\^{}.p)/(16*(d - 1))} & \texttt{-p.p\^{}2*(d*a.b*c\^{}.p - 3*a.b*c\^{}.p + a.c*b\^{}.p + a\^{}.p*b.c - a\^{}.p*b~.c~ + 2*a~.b~*c\^{}.p - a~.c~*b\^{}.p)/(16*(d - 1))} & \texttt{0} & \texttt{0} \\
\texttt{QdS (hat)} & \texttt{p.p\^{}2*(d*a\^{}.b\^{}*c\^{}.p - 3*a\^{}.b\^{}*c\^{}.p + a\^{}.c\^{}*b\^{}.p + a\^{}.p*b\^{}.c\^{})/(16*(d - 1))} & \texttt{-p.p\^{}2*(d*a\^{}.b\^{}*c\^{}.p - 3*a\^{}.b\^{}*c\^{}.p + a\^{}.c\^{}*b\^{}.p + a\^{}.p*b\^{}.c\^{})/(16*(d - 1))} & \texttt{0} & \texttt{0} \\
\texttt{c tilde, h hat} & \texttt{0} & \texttt{0} & \texttt{0} & \texttt{0} \\
\texttt{c tilde, h mixed} & \texttt{0} & \texttt{0} & \texttt{0} & \texttt{0} \\
\texttt{c hat, h mixed} & \texttt{0} & \texttt{0} & \texttt{0} & \texttt{0} \\
\texttt{all tilde} & \texttt{0} & \texttt{0} & \texttt{0} & \texttt{0} \\
\texttt{c hat, h tilde} & \texttt{a~.b~*c\^{}.p*p.p\^{}2/16} & \texttt{-a~.b~*c\^{}.p*p.p\^{}2/16} & \texttt{0} & \texttt{0} \\
\hline\end{tabular}\end{center}
\noindent Full-row verdict: \texttt{DECOUPLED (proof, full row): p\_m G = 0 exactly in d for D11 by itself and for D01 + D03 as a sum; no evanescent structure survives in the slot}.

